\chapter{Các luồng nghiệp vụ}
\label{ch:luong-nghiep-vu}

Nghiệp vụ của hệ thống gồm năm luồng. 
Hình~\ref{fig:activity-overview} và hình~\ref{fig:vong-hoc} mô tả cùng một luồng, nhưng ở hai mức chi tiết khác nhau. 

\begin{landscape}
\begin{figure}[p]
  \centering
  \includegraphics[width=\linewidth,height=\textheight,keepaspectratio]{activity-overview.pdf}
  \caption[Activity Overview]{Activity Overview --- toàn bộ vòng, hai giai đoạn nằm trong hai khung
  riêng}
  \label{fig:activity-overview}
\end{figure}
\end{landscape}

\section{Workflow 1 --- Tạo và duyệt đề}

\paragraph{Mục tiêu.} Giúp giáo viên tạo được đề phù hợp với mục tiêu học tập, nhưng giáo viên
vẫn là người duyệt cuối cùng trước khi phát hành.

\paragraph{Luồng chính.}
\begin{enumerate}
  \item Teacher nhập yêu cầu tạo đề.
  \item Teacher cung cấp ràng buộc: chủ đề, độ khó, số lượng câu hỏi, loại bài kiểm tra, thời
        gian làm bài.
  \item \texttt{System/AI} phân tích yêu cầu.
  \item \texttt{System/AI} tạo đề nháp gồm câu hỏi, phương án trả lời, đáp án đúng, lời giải và
        metadata cơ bản.
  \item Teacher xem đề nháp.
  \item Teacher chỉnh sửa thủ công, hoặc yêu cầu tạo lại một phần.
  \item Teacher duyệt đề.
  \item Teacher phát hành đề cho Student.
\end{enumerate}

\paragraph{Điểm kiểm soát.} Teacher phải duyệt đề trước khi đề được phát hành. Đề kiểm tra ảnh
hưởng trực tiếp đến việc đánh giá học sinh, nên trợ lí không được tự phát hành.

\section{Workflow 2 --- Làm bài (giai đoạn 1)}

\paragraph{Mục tiêu.} Cho phép học sinh làm bài và ghi nhận bài nộp để hệ thống đánh giá.

\paragraph{Luồng chính.}
\begin{enumerate}
  \item Student nhận assessment đã được Teacher phát hành.
  \item Student vào làm, \emph{tới hết} giờ đóng.
  \item Student trả lời câu hỏi trắc nghiệm.
  \item Student nộp bài.
  \item Hệ thống ghi nhận submission để đánh giá.
\end{enumerate}

\subsection{Ba mốc thời gian}
\label{sec:ba-con-so}

Giai đoạn 1 có ba tham số thời gian độc lập: thời gian làm bài, giờ mở và giờ đóng.

Luật là: giờ đóng chặn việc \emph{vào làm}, không chặn việc nộp. Học sinh vào tới hết giờ đóng vẫn
được. Ai đã vào rồi thì không bị dừng giữa chừng, và được đủ thời gian làm bài kể cả khi tràn qua giờ đóng.
Ví dụ, đóng lúc 18:00 với thời gian làm bài 15 phút thì bài cuối có thể nộp lúc 18:15.

\section{Workflow 3 --- Chấm bài và sinh nhận xét}

\paragraph{Mục tiêu.} Đánh giá bài làm, xác định lỗi sai hoặc hiểu nhầm, và đưa nhận xét phù hợp.

\paragraph{Luồng chính.}
\begin{enumerate}
  \item \texttt{System/AI} nhận submission.
  \item So sánh với đáp án chuẩn và lời giải chuẩn.
  \item Hệ thống \textbf{tra} lỗi sai từ phương án nhiễu đã chọn --- mỗi nhiễu được soạn kèm một lỗi cụ thể.
  \item Tạo kết quả gồm điểm, nhận xét, lỗi sai, mã hiểu nhầm và độ tin cậy.
  \item Hệ thống quyết định kết quả có cần giáo viên xem lại hay không.
\end{enumerate}

\paragraph{Nguyên tắc đánh giá.}
\begin{itemize}
  \item Đáp án đúng nhưng cách làm sai \emph{không} nên mặc định là đã hiểu.
  \item Đáp án sai nhưng một phần cách làm đúng cần được ghi nhận.
  \item Mỗi phương án nhiễu phải đại diện cho một lỗi cụ thể, soạn cùng câu hỏi. Đây là điều kiện
        để giai đoạn 2 giải thích được.
  \item Nhận xét nên giúp học sinh hiểu lỗi và biết bước tiếp theo.
\end{itemize}

\section{Workflow 4 --- Hàng đợi review}

\paragraph{Mục tiêu.} Đưa các kết quả có độ tin cậy thấp hoặc có dấu hiệu bất thường cho giáo
viên xem xét.

\paragraph{Luồng chính.}
\begin{enumerate}
  \item \texttt{System/AI} hoàn thành đánh giá và tính độ tin cậy ở mức nghiệp vụ.
  \item Độ tin cậy đủ cao thì kết quả đi thẳng sang giai đoạn 2.
  \item Độ tin cậy thấp thì kết quả đi vào Teacher Review Queue.
  \item Teacher xem câu hỏi, đáp án chuẩn, bài làm, nhận xét của trợ lí, và lý do cần xem lại.
  \item Teacher xác nhận hoặc chỉnh sửa điểm, lỗi sai, hiểu nhầm và nhận xét.
  \item Hệ thống lưu kết quả cuối cùng.
\end{enumerate}

\paragraph{Nhánh này không chặn giai đoạn 2.} Đây là điểm đáng chú ý: một kết quả đi vào hàng đợi
\emph{không} giữ học sinh lại. Teacher chốt xong thì luồng hợp lại, còn học sinh vẫn chữa bài
bình thường trong lúc chờ.

\paragraph{Teacher bắt buộc xem lại khi:} độ tin cậy thấp; bài làm có mâu thuẫn, ví dụ đáp án
đúng nhưng cách làm sai; hệ thống không đủ căn cứ để xác định lỗi; hoặc có trường hợp bất thường cần
người có chuyên môn xem lại.

\section{Workflow 5 --- Chữa bài (giai đoạn 2)}

\paragraph{Mục tiêu.} Học sinh sửa được lỗi vừa mắc, và bài kiểm tra chỉ kết thúc khi việc đó đã
diễn ra.

\paragraph{Luồng chính.}
\begin{enumerate}
  \item Học sinh trò chuyện với trợ lí Kriky về lỗi sai, chỗ nào chưa hiểu, cách giải như thế nào. 
  \item Kriky dựa vào \emph{lời giải đã soạn kèm câu hỏi}. Học sinh hỏi lại đến khi hiểu. Phần này \textbf{không tính giờ}.
  \item Trong lúc đó, hệ thống sinh câu biến thể của câu mà học sinh làm sai.
  \item Học sinh bấm nút làm bài mới. Đồng hồ của lượt bắt đầu chạy.
  \item Học sinh làm các câu mới trong lượt.
  \item Câu nào đúng thì chốt \textbf{0,5} điểm; câu nào còn sai thì sang vòng tiếp.
  \item Mỗi câu tối đa \textbf{ba vòng}. Hết ba vòng mà vẫn sai thì chốt \textbf{0} điểm.
  \item Bài kết thúc khi mọi câu đã chốt, hoặc khi hết hạn giai đoạn 2.
\end{enumerate}

\subsection{Giai đoạn 2 đếm giờ như thế nào?}
\label{sec:hai-dong-ho}

Giai đoạn 2 nhận hai tham số: số phút mỗi câu, và hạn kết thúc giai đoạn 2.

Thời lượng của một lượt bằng số phút mỗi câu nhân với số câu còn sai trong lượt đó, nên sai hai câu với 
thời lượng làm lại 5 phút mỗi câu thì lượt ấy dài 10 phút.

Phần giải thích không tính giờ. Đồng hồ chỉ bắt đầu khi học sinh tự bấm nút làm bài mới.

\subsection{Khi nào kết thúc giai đoạn 2}

Hạn kết thúc giai đoạn 2 là một mốc tuyệt đối do giáo viên đặt, chung cho cả lớp, và dài hơn hẳn hạn
của giai đoạn 1, chẳng hạn là cuối ngày.

Hết hạn thì dừng, kể cả khi học sinh đang làm dở một lượt, và câu chưa xong nhận 0 điểm. 

Đây là chỗ giai đoạn 2 làm ngược với giai đoạn 1. Ở giai đoạn 1, học sinh đã vào thì không bị dừng
giữa chừng; ở giai đoạn 2, hết hạn là dừng. Hai luật nghịch nhau một cách có chủ ý, và
chương~\ref{ch:quyet-dinh} giải thích vì sao.

\subsection{Giáo viên làm gì?}

Học sinh báo cáo một câu hoặc một đoạn hội thoại là \emph{giải thích chưa rõ}. Giáo viên xem được các
báo cáo này và giải thích để đưa vào bài dạy trên lớp.

\section{Business Workflows}

Hình~\ref{fig:activity-overview} cho thấy luồng tổng quan;
hình~\ref{fig:business-workflows} cho thấy trách nhiệm của từng vị trí.

\begin{figure}[p]
  \centering
  \includegraphics[width=\linewidth,height=\textheight,keepaspectratio]{business-workflows.pdf}
  \caption[Business Workflows]{Business Workflows --- cùng vòng ấy, tách theo lane}
  \label{fig:business-workflows}
\end{figure}